\documentclass[11pt]{article}
\usepackage[margin=1in]{geometry}
\usepackage{amsmath,amssymb,amsthm}

\newcommand{\ALCIRCC}{\mathcal{ALCI}_{\mathrm{RCC5}}}
\newcommand{\DR}{\mathsf{DR}}
\newcommand{\PO}{\mathsf{PO}}
\newcommand{\PP}{\mathsf{PP}}
\newcommand{\PPI}{\mathsf{PPI}}
\newcommand{\EQ}{\mathsf{EQ}}
\newcommand{\Comp}{\mathrm{comp}}
\newcommand{\cl}{\mathrm{cl}}
\newcommand{\code}[1]{\texttt{#1}}

\theoremstyle{plain}
\newtheorem{theorem}{Theorem}
\newtheorem{lemma}[theorem]{Lemma}
\newtheorem{corollary}[theorem]{Corollary}
\theoremstyle{definition}
\newtheorem{definition}[theorem]{Definition}
\newtheorem{remark}[theorem]{Remark}

\title{A decidable fragment of $\ALCIRCC$:\\ the $\forall\PO$-free
concepts}
\author{}
\date{2026-07-16}

\begin{document}
\maketitle

\begin{abstract}
Concept satisfiability for $\ALCIRCC$ --- the description logic
$\mathcal{ALCI}$ over the five RCC5 base relations under abstract
composition-table semantics --- is open (Wessel 2002/2003), and reduces
to a bounded-live-width property F6 whose status is the standing
obstruction. We isolate a syntactically defined, unconditionally
decidable fragment: concepts with \emph{no $\forall\PO$ subformula}. The
key algebraic fact is that $\PO$ is the only relation whose recurrence
along a vertical chain requires an all-phase safety condition not
supplied by backward forcing or forward absorption; absent
$\forall\PO$, that condition is vacuous, the two-tier quotient is
finite, and the live width is bounded by a computable function of the
concept. The fragment retains $\exists\PO$, $\forall\DR$, $\exists\DR$
and arbitrary vertical ($\PP/\PPI$) modalities. This note is a
theorem-level write-up; its proof rests on the lemmas of the two-tier
quotient development (\code{papers/two\_tier\_quotient\_ALCIRCC5.tex})
and is not yet machine-certified. Its provenance is a cold research
attack on the F6/W2$'$ problem (GPT-5.5); we record it here with that
status.
\end{abstract}

\section{The fragment}

We work in negation normal form (NNF) with inverse roles absorbed
($\exists\PP^-.D = \exists\PPI.D$), under strong-EQ semantics ($\EQ$ is
identity; between distinct elements only $\{\DR,\PO,\PP,\PPI\}$ occur),
and abstract composition-table semantics: a model is a set with a
total, converse-closed, composition-closed labelling of ordered pairs by
RCC5 relations. Let $\cl(C)$ be the closure of $C$ (subformulas closed
under the usual Fischer--Ladner conditions) and $m=|\cl(C)|$.

\begin{definition}[$\forall\PO$-free]
A concept $C$ is \emph{$\forall\PO$-free} if no subformula of $C$ has the
form $\forall\PO.D$.
\end{definition}

The fragment is strictly weaker than banning ``horizontal universals
under horizontal existentials'': it still allows $\exists\PO.D$,
$\forall\DR.D$, $\exists\DR.D$, and unbounded vertical chains, so it
retains infinite models and genuine horizontal existentials.

\section{Why $\PO$ is the unstable case}

Recall the four determining (singleton) compositions of proper
relations,
\[
\Comp(\DR,\PPI)=\{\DR\},\quad
\Comp(\PP,\DR)=\{\DR\},
\]
\[
\Comp(\PP,\PP)=\{\PP\},\quad
\Comp(\PPI,\PPI)=\{\PPI\},
\]
and the two relevant non-singleton facts
\[
\Comp(\PP,\PO)=\{\DR,\PO,\PP\},\qquad
\Comp(\PPI,\PO)=\{\PO,\PPI\}.
\]
Along a vertical chain $d_0\mathbin{\PP}d_1\mathbin{\PP}d_2\cdots$, the
relation of an external witness $e$ to the chain \emph{stabilizes}: the
RCC5 transition table admits only finitely many monotone exits among
$\DR,\PO,\PP,\PPI$ (external-relation stabilization). For the three
``stable'' relations this gives an all-phase witness by a determining
composition:
\begin{itemize}
\item a $\DR$-demand is carried to all deeper phases by
$\Comp(\PP,\DR)=\{\DR\}$ (backward forcing);
\item a $\PP$-demand by $\Comp(\PP,\PP)=\{\PP\}$;
\item a $\PPI$-demand to all shallower phases by
$\Comp(\PPI,\PPI)=\{\PPI\}$ (forward absorption).
\end{itemize}
$\PO$ has \emph{neither}: $\Comp(\PP,\PO)$ and $\Comp(\PPI,\PO)$ are
non-singletons, so a witness that is $\PO$-related at one phase need not
be $\PO$-related at another. This is precisely the point at which the
two-tier quotient's PO-coherence condition can be violated --- but only
by a \emph{universal} $\PO$-obligation.

\begin{lemma}[Automatic PO-coherence]\label{lem:pocoh}
Every $\forall\PO$-free concept is PO-coherent in the period-descriptor
sense: for every recurrent descriptor $\mathrm{Desc}=(\tau_1,\dots,
\tau_p)$ and every $\exists\PO.D\in\bigcup_j\tau_j$, any Hintikka type
$\sigma$ with $D\in\sigma$ is $\PO$-safe for all phases $\tau_j$ with
respect to $\PO$-universal obligations.
\end{lemma}

\begin{proof}
The type-safety test for a $\PO$-edge $\tau_j\to\sigma$ requires, for the
$\PO$-universal part, that if $\forall\PO.E\in\tau_j$ then $E\in\sigma$,
and (since $\PO$ is symmetric) the analogous condition from $\sigma$ to
$\tau_j$. In a $\forall\PO$-free concept no type contains any
$\forall\PO.E$, so the $\PO$-universal part of $\mathrm{Safe}(\tau_j,
\sigma)$ is empty for every phase. The remaining requirements are the
propositional/type conditions already secured by $D\in\sigma$; hence the
same $\sigma$ is safe for all phases.
\end{proof}

\section{Bounded quotient and decidability}

\begin{theorem}[Bounded live width; decidability]\label{thm:main}
For every satisfiable $\forall\PO$-free concept $C$ there is a
split-forest presentation of a model of $C$ with live width at most
\[
K(C)=(1+m)^2\,2^{2^m},\qquad m=|\cl(C)|.
\]
Consequently concept satisfiability for the $\forall\PO$-free fragment
of $\ALCIRCC$ is decidable.
\end{theorem}

\begin{proof}[Proof sketch]
Start from any model of $C$. Along each infinite $\PP$-chain record the
recurrent set of Hintikka types with its cyclic order as a
\emph{period descriptor}; there are at most $2^m$ types and hence at most
$D(C)\le 2^{2^m}$ descriptors. By external-relation stabilization,
$\DR/\PP$-demands obtain all-phase witnesses through backward forcing,
$\PPI$-demands through forward absorption, and (Lemma~\ref{lem:pocoh})
$\PO$-demands obtain all-phase safety because there are no $\forall\PO$
constraints. Thus every recurrent descriptor has finitely many all-phase
witnesses, at most one regular/designated witness per existential demand
per descriptor, giving a two-tier quotient of size
$Q(C)\le D(C)(1+m)\le(1+m)\,2^{2^m}$.

The quotient network is checked by ordinary RCC5 path-consistency, which
the patchwork property (Renz--Nebel) turns into global realizability;
unfolding each kernel as its periodic $\PP$-chain yields a model of $C$.
All repeated vertical material is folded into the kernel loop, and
cross-edges to external witnesses are constant quotient edges or
vertical shadows, so the live interface width is bounded by the quotient
size times the local number of existential demands:
$K(C)\le Q(C)(1+m)$. Decidability follows by enumerating all quotients
up to this computable bound and checking the finite type, witness,
safety, and RCC5 composition conditions.
\end{proof}

\begin{remark}[What the bound is and is not]
The constant is not meant to be tight; only its \emph{computability}
matters, as decidability then follows from the conditional theorem of
the width-barrier development (a computable live-width bound inhabits F6
for the fragment). The exponents reflect the type/descriptor counting,
not a claimed complexity result.
\end{remark}

\begin{remark}[Relation to W2$'$]
The mathematical fragment result uses an explicit quotient with
designated witnesses; it does not require arbitrary glue steps to factor
through a fixed coarse interface state, so it does not depend on the
(false) coarse uniformization property W2$'$. A \emph{catalogue}
implementation with a fixed finite steering syntax would still need
row-coloured steering or a principal-row-profile refinement, but the
decidability of the fragment does not.
\end{remark}

\begin{remark}[Dependence on the two-tier Step~2, repaired 2026-08-28]
\label{rem:step2dep}
The quotient above needs one relation per ordered pair of kernels, which the
two-tier development supplies in Step~2 of its completeness theorem. That step
originally defined it as a double limit $\lim_{i,j\to\infty}\rho(d_i,d_j)$,
justified by external-relation stabilization ``applied twice''. The inference is
invalid --- the stabilization lemma is \emph{one-sided}, fixing an element and
varying the chain, and the double limit need not exist --- so the step was a
genuine gap. This paper never states it: it uses stabilization one-sidedly, as
the lemma allows. But it builds a two-tier quotient, so it inherits the step,
and the reader should know where the repair is. It is in
\code{papers/two\_tier\_quotient\_ALCIRCC5.tex} (Remark ``The lemma is
one-sided'' and Lemma ``Simultaneous finite rectangles''), which replaces the
limit by a constant finite \emph{rectangle} --- all that a quotient reading
finitely many phases needs --- and it is machine-checked in Lean as
\code{fused\_kq\_all}. Theorem~\ref{thm:main} and the bound $K(C)$ are
unaffected.
\end{remark}

\section{Provenance and status}

This fragment and its proof were produced by a cold research attack on
the F6/W2$'$ problem (GPT-5.5, on the packet
\code{papers/cold\_review\_f6\_w2prime}; the response and its supporting
scripts are archived under \code{attack\_req/}). The underlying
ingredients --- external-relation stabilization, backward-forcing safety
for $\DR/\PP$, forward absorption for $\PPI$, and the PO-coherence
condition --- are those of the two-tier quotient development
(\code{papers/two\_tier\_quotient\_ALCIRCC5.tex}), specialized here by a
syntactic sufficient condition (Lemma~\ref{lem:pocoh}). The result is
\emph{theorem-level and not machine-certified}; it should be read as a
proof obligation whose lemmas are already present in that development,
and it should be independently reviewed (and ideally Lean-transcribed)
before being cited as established. Consistent with the project's standing
discipline, the honest label remains: \emph{strongly supported, not
certified}.

\section{Relationship to the Lean artifact (added 2026-08-20)}
\label{sec:leanrel}

Since this paper was written, a substantial Lean development for this same
fragment has grown in \code{formal/POFreeLift.lean}. The relationship is
\emph{not} that the Lean transcribes the proof above, and the two must not be
conflated when citing either. Precisely:

\paragraph{Shared engine.} Every ingredient named in the sketch of
Theorem~\ref{thm:main} is machine-checked in the Lean file, under these names:
external-relation stabilization (\code{external\_stabilizes}), backward forcing
for $\DR$ and $\PP$ (\code{backward\_forcing\_dr}, \code{backward\_forcing\_pp}),
forward absorption for $\PPI$ (\code{forward\_absorption\_ppi}), and the
$\forall\PO$-vacuity that makes Lemma~\ref{lem:pocoh} free
(\code{mty\_no\_all\_po}, resting on \code{pofree\_cl\_all}).

\paragraph{Different final step, and what is therefore NOT certified.} This paper
concludes with a \emph{split-forest presentation} of live width at most
$K(C)=(1+m)^2 2^{2^m}$, and derives decidability from that bound. The Lean
development instead builds a finite \emph{certificate} (\code{MultiTier}:
$\kappa$ periodic kernels plus $\beta$ externals) and obtains
\code{Decidable (Satisfiable C0)} from a soundness theorem
(\code{multiTier\_sound}) together with finite enumeration of certificate codes.
Consequently the Lean file contains \emph{no width statement at all}: neither
$K(C)$, nor live width, nor the existence of a split-forest presentation, is
machine-checked. A reader must not take the Lean development as certifying
Theorem~\ref{thm:main} as stated.

\paragraph{The split-forest idea is nonetheless load-bearing.} The certificate's
shape \emph{is} the split-forest discipline: kernels are the vertical
$\PP/\PPI$ skeleton (with a direction flag), externals carry the horizontal
$\DR/\PO$ completion, and the interface rows are the cross edges. So the
decomposition remains the semantic foundation, exactly as the overview claims;
what the Lean route dispenses with is the split forest as a formal object
carrying a width measure, not the insight.

\paragraph{Consequence for a rewrite.} Three registers are easy to conflate and
should be separated explicitly in any revision: (i) the split-forest
\emph{insight}, which is alive and formalized in structure; (ii) the
split-forest \emph{presentation with its $K(C)$ bound}, which is this paper's
theorem and is a proof sketch, unreviewed and not machine-checked; (iii) the
split-forest \emph{plus patchwork bags and a parity automaton} of the
full-logic argument (Theorems A/B/C), which is untouched by the Lean work and
still gated on F6.

\end{document}
